\SourceLeaf{350}%
relations must be distinguished. Some presuppose a definite shape of the surface in space. Others are independent of the different shapes the surface can take. The latter are the subject here. As the preceding discussion shows, the measure of curvature belongs to this class. The same is true, as one can readily see, of the study of figures drawn on the surface, their angles, areas and total curvature, of joining points by shortest lines, and of similar questions. All such investigations must start from the premise that the nature of the curved surface itself is given by the expression for an arbitrary line element of the form \(\BareRoot(E\mathrm{d}p^2+2F\mathrm{d}p\mathord{.}\mathrm{d}q+G\mathrm{d}q^2)\). In the present memoir the author has included some of the investigations he has pursued in this field for several years. He has limited himself to those that remain close to the starting point and some of which can serve as general tools for many further investigations. In our notice we must narrow the scope further and give just a few examples. The following theorems may serve.

If infinitely many shortest lines of equal length start from a single point on a curved surface, the line through their endpoints intersects each of them at right angles. If shortest lines of equal length are drawn from each point of an arbitrary line on a curved surface, perpendicular to that line, then they are also all perpendicular to the line joining their other endpoints. Both theorems, the second of which can be viewed as a generalization of the first, are proved analytically and by simple geometric arguments. \textit{The sum of the angles of a triangle formed by shortest lines exceeds two right angles by the total curvature of the triangle.} The unit for angles is taken to be the angle subtended by an arc equal to the radius \((57^0\ 17'\ 45'')\). For total curvature as a patch of the auxiliary sphere, the unit is the area of a square whose side equals the radius of that sphere. This important theorem can clearly be stated another way. The excess of the angles of a triangle formed by shortest lines over two right angles is to eight right angles as the corresponding patch of the surface of the %
